\chapter{Il background} % 2

\section{Malware} % 2.1
\subsection{Botnet} % 2.1.1
\subsubsection{Botnet IoT}
\subsection{Command-and-Control} % 2.1.2
\subsubsection{Handshake, profiling e beaconing}
\subsubsection{Attacchi}


\section{Malware Analysis} % 2.2

Fonti Bibliografiche: Practical Malware Analysis: The Hands-On Guide to Dissecting Malicious Software

Argomenti da trattare:

\begin{itemize}
    \item Definizione di malware analysis
    \item Obiettivi della malware analysis
    \item Breve introduzione/conclusione
\end{itemize}

\subsection{Analisi statica} % 2.2.1

Fonti Bibliografiche: Practical Malware Analysis: The Hands-On Guide to Dissecting Malicious Software; Limits of Static Analysis for Malware Detection

Argomenti da trattare:

\begin{itemize}
    \item Tecniche di analisi statica base e avanzate
    \item Efficacia e limitazioni dell'analisi statica
    \item Rimozione della tabella dei simboli / stripping del malware (fonte?)
\end{itemize}

\subsection{Analisi dinamica} % 2.2.2

Fonti Bibliografiche: Practical Malware Analysis: The Hands-On Guide to Dissecting Malicious Software; A Survey on Automated Dynamic Malware-Analysis Techniques and Tools

Argomenti da trattare:

\begin{itemize}
    \item Tecniche di analisi dinamica base e avanzate
    \item Ostacoli e rischi dell'analisi dinamica (fonte?)
\end{itemize}

\subsection{Disassembly e decompiling} % 2.2.3

Fonti Bibliografiche: Practical Malware Analysis: The Hands-On Guide to Dissecting Malicious Software; The Ghidra Book: The Definitive Guide; https://www.ghidradocs.com/12.1.2\_PUBLIC/help/index.html; The IDA Pro Book; https://docs.hex-rays.com/

Argomenti da trattare:

\begin{itemize}
    \item Definizione e differenza tra disassembly e decompiling
    \item Strumenti di reversing, Ghidra e IDA
\end{itemize}

\subsection{Reverse engineering} % 2.2.4

Fonti Bibliografiche: ???

Argomenti da trattare:

\begin{itemize}
    \item Definizione Generica?
\end{itemize}

\subsubsection{Assistito}

Fonti Bibliografiche: LLM4Decompile: Decompiling Binary Code with Large Language Models; ReCopilot: Reverse Engineering Copilot in Binary Analysis; Decompiling the Synergy: An Empirical Study of Human--LLM Teaming in Software Reverse Engineering

Argomenti da trattare:

\begin{itemize}
    \item Esempi di come il reverse engineering può essere assistito dai large language models
\end{itemize}

\section{Large Language Models} % 2.3

Fonti Bibliografiche: ???

Argomenti da trattare:

\begin{itemize}
    \item Generative AI
    \item Large Language Models
\end{itemize}

\subsection{Prompt Engineering} % 2.3.1

Fonti Bibliografiche: ???

Argomenti da trattare:

\begin{itemize}
    \item Prompt engineering (in breve)
\end{itemize}

\subsection{AI Agents e Tools} % 2.3.2

Fonti Bibliografiche: ???, ReAct: Synergizing Reasoning and Acting in Language Models, ???

Argomenti da trattare:

\begin{itemize}
    \item Reasoning e chain of thought
    \item Intelligent Agents
    \item Long horizon / one-shot tasks (?)
    \item Tools e tool calls
\end{itemize}

\subsection{Model Context Protocol} % 2.3.3

Fonti Bibliografiche: What is the Model Context Protocol (MCP)

Argomenti da trattare:

\begin{itemize}
    \item Cos'è il protocollo MCP (obiettivo)
    \item Architettura di host, server e client.
\end{itemize}

% long horizon one shots ai agents + chain of thought (native reasoning)
% https://arxiv.org/html/2509.09677v3